\documentclass[11pt,a4paper]{article}
\newcommand{\fichetype}{Corrigé détaillé}
% Thème A4 commun, librement modifiable.
\usepackage[T1]{fontenc}
\usepackage[utf8]{inputenc}
\usepackage[french]{babel}
\usepackage{lmodern,amsmath,amssymb,mathtools,tikz,pgfplots}
\usepackage[margin=2.2cm,headheight=15pt]{geometry}
\usepackage{xcolor,fancyhdr,enumitem,hyperref,needspace}
\usetikzlibrary{arrows.meta,positioning,calc}
\pgfplotsset{compat=1.18,width=11cm,height=6cm}
\definecolor{Ink}{HTML}{202124}
\definecolor{Blue}{HTML}{356B9E}
\definecolor{Muted}{HTML}{737579}
\color{Ink}
\renewcommand{\familydefault}{\sfdefault}
\setlength{\parindent}{0pt}
\setlength{\parskip}{6pt}
\setlist{itemsep=3pt,topsep=4pt}
\pagestyle{fancy}\fancyhf{}
\fancyhead[L]{\small\color{Muted}MATHS LYCÉE}
\fancyhead[R]{\small\color{Muted}2026--2027}
\fancyfoot[L]{\small\color{Muted}\fichetype}
\fancyfoot[R]{\small\thepage}
\renewcommand{\headrulewidth}{0.3pt}
\hypersetup{colorlinks=true,urlcolor=Blue,linkcolor=Blue}
\newcommand{\R}{\mathbb{R}}
\newcommand{\N}{\mathbb{N}}
\newcommand{\E}{\mathbb{E}}
\newcommand{\V}{\mathrm{V}}
\newcommand{\exercice}[3]{\Needspace{6\baselineskip}\vspace{0.5cm}{\color{Blue}\small\textbf{EXERCICE #1}\hfill #2}\par{\large\bfseries #3}\par}

\begin{document}
{\small\color{Blue}PREMIÈRE SPÉCIALITÉ}\par{\LARGE\bfseries Nombre dérivé, tangente et approximation}\par{\large Corrigé détaillé}\par\vspace{0.2cm}\hrule\vspace{0.4cm}
\par\noindent\begin{minipage}{\linewidth}\exercice{1}{Application}{Taux moyens}
{\small\color{Muted}\textbf{Énoncé.} Pour $f(x)=2x^2+1$, calculer les taux entre $1$ et $3$, puis entre $1$ et $1,1$. Interpréter géométriquement.}
\par\medskip
\textbf{Résolution.}\par
$f(1)=3$ et $f(3)=19$, donc le premier taux est $(19-3)/2=8$. Puis $f(1,1)=3,42$, donc le second vaut $(3,42-3)/0,1=4,2$. Ce sont les pentes des sécantes passant par les deux points correspondants de la courbe.
\end{minipage}\par
\par\noindent\begin{minipage}{\linewidth}\exercice{2}{Entraînement}{Une dérivée par définition}
{\small\color{Muted}\textbf{Énoncé.} Pour $f(x)=3x^2-2x$, calculer le taux en $a$ avec l’accroissement $h\ne0$, puis $f\prime(a)$.}
\par\medskip
\textbf{Résolution.}\par
On développe $f(a+h)=3a^2+6ah+3h^2-2a-2h$. En soustrayant $f(a)$ puis en divisant par $h$, on obtient $6a-2+3h$. Quand $h$ tend vers $0$, ce taux tend vers $6a-2$, donc $f\prime(a)=6a-2$.
\end{minipage}\par
\par\noindent\begin{minipage}{\linewidth}\exercice{3}{Application}{Équation de tangente}
{\small\color{Muted}\textbf{Énoncé.} $f(x)=x^2-4x+1$. On admet $f\prime(x)=2x-4$. Déterminer la tangente en $a=3$ et son intersection avec l’axe des ordonnées.}
\par\medskip
\textbf{Résolution.}\par
$f(3)=9-12+1=-2$ et $f\prime(3)=2$. L’équation est $y=-2+2(x-3)=2x-8$. À l’axe des ordonnées, $x=0$, donc le point d’intersection est $(0,-8)$. Vérifier $x=3$ redonne $y=-2$.
\end{minipage}\par
\par\noindent\begin{minipage}{\linewidth}\exercice{4}{Application}{Lecture d’une tangente}
{\small\color{Muted}\textbf{Énoncé.} La courbe passe par $A(-1,4)$. Sa tangente en $A$ passe par $B(2,-2)$. Donner $f(-1)$, $f\prime(-1)$ et l’équation de cette tangente.}
\par\medskip
\textbf{Résolution.}\par
Le point $A$ donne $f(-1)=4$. La pente de $AB$ vaut $(-2-4)/(2-(-1))=-6/3=-2$, donc $f\prime(-1)=-2$. L’équation est $y=4-2(x+1)=-2x+2$. Le point $B$ est un point de la tangente ; on ne peut pas en déduire $f(2)=-2$.
\end{minipage}\par
\par\noindent\begin{minipage}{\linewidth}\exercice{5}{Entraînement}{Approximer un carré}
{\small\color{Muted}\textbf{Énoncé.} Utiliser la tangente à la fonction carré en $5$ pour approcher $5,03^2$. Calculer ensuite l’erreur exacte et son signe.}
\par\medskip
\textbf{Résolution.}\par
Pour $f(x)=x^2$, $f(5)=25$ et $f\prime(5)=10$. Avec $h=0,03$, l’approximation vaut $25+10\times0,03=25,3$. Le développement exact donne $(5+0,03)^2=25,3+0,0009=25,3009$. L’approximation est inférieure de $0,0009$ à la valeur exacte.
\end{minipage}\par
\par\noindent\begin{minipage}{\linewidth}\exercice{6}{Entraînement}{Une vitesse instantanée}
{\small\color{Muted}\textbf{Énoncé.} Une distance fictive vaut $d(t)=2t^2$ mètres. Calculer la vitesse moyenne entre $3$ et $3+h$, puis la vitesse instantanée en $3$.}
\par\medskip
\textbf{Résolution.}\par
Le taux est $[2(3+h)^2-18]/h=(12h+2h^2)/h=12+2h$ m/s, pour $h\ne0$ et des temps admissibles. Il tend vers $12$ quand $h$ tend vers $0$. La vitesse instantanée en $3$ vaut donc $12$ m/s. Il s’agit d’une vitesse, pas de la distance $d(3)=18$ m.
\end{minipage}\par
\par\noindent\begin{minipage}{\linewidth}\exercice{7}{Entraînement}{Une liste de sécantes}
{\small\color{Muted}\textbf{Énoncé.} Pour $f(x)=1/x$, calculer les taux en $1$ pour $h=0,1$ et $h=-0,1$. Décrire un algorithme produisant les taux pour $h=10^{-k}$, $k=1,\ldots,5$.}
\par\medskip
\textbf{Résolution.}\par
Le taux simplifié est $-1/(1+h)$. Pour $h=0,1$, il vaut $-10/11\approx-0,9091$ ; pour $h=-0,1$, il vaut $-10/9\approx-1,1111$. Pour chaque entier $k$ de $1$ à $5$, poser $h=10^{-k}$ et ajouter $-1/(1+h)$ à une liste. Les valeurs s’approchent de $-1=f\prime(1)$. La forme simplifiée évite la soustraction de nombres presque égaux.
\end{minipage}\par
\par\noindent\begin{minipage}{\linewidth}\exercice{8}{Approfondissement}{Deux fonctions en zéro}
{\small\color{Muted}\textbf{Énoncé.} Comparer la dérivabilité en $0$ de $f(x)=|x|$ et $g(x)=x|x|$. Utiliser les quotients de variation.}
\par\medskip
\textbf{Résolution.}\par
Pour $f$, le quotient est $|h|/h$, égal à $1$ à droite et $-1$ à gauche : pas de dérivabilité. Pour $g$, $g(0)=0$ et le quotient est $h|h|/h=|h|$, qui tend vers $0$ des deux côtés. Donc $g$ est dérivable en $0$ et $g\prime(0)=0$. La présence d’une valeur absolue dans une formule ne suffit pas à conclure à la non-dérivabilité.
\end{minipage}\par
\end{document}
